\begin{EMtheorem}{The residue theorem}

If the function \EMim{f(z)} is analytic everywhere inside the path \EMim{C} except at the points \EMim{z_1,z_2,\dots,z_n}, and the residues of the function at these points are \EMim{A_1^1,A_1^2,\dots,A_1^n} respectively, then
\EMdisp{\oint_Cf(z)\,dz\EMbr =2\pi i\sum_{k=1}^{n}A_1^k}

\end{EMtheorem}

Many different kinds of integrals can be solved by the method of residues. Two classes of them are mentioned here:

\begin{EMdefinition}{Integrals of the first class}

In these integrals the denominator of the integrand is a polynomial, and the numerator is a polynomial, an exponential, or a sine or cosine function (which can also be expressed exponentially).

\end{EMdefinition}

\begin{EMdefinition}{Integrals of the second class}

In these integrals the integrand is a function of sine or cosine.

\end{EMdefinition}

\EMhold

\EMsection{Integrals of the First Class}{Integrals of the First Class}

\begin{EMexample}{}

Evaluate the integral \EMim{\displaystyle\int_0^\infty\frac{\cos\omega x}{(1+x^2)^2}\,dx} (\EMim{\omega>0}).

Consider the function \EMim{f(z)=\dfrac{e^{i\omega z}}{(1+z^2)^2}}. It is not analytic at \EMim{z=i} and \EMim{z=-i}. The path \EMim{C} is the upper semicircle of radius \EMim{R}.
\EMdisp{\begin{aligned}
\oint_C\frac{e^{i\omega z}}{(1+z^2)^2}\,dz&=2\pi i\times\operatorname{Res}\left\{\frac{e^{i\omega z}}{(1+z^2)^2}\right\}\Bigg|_{z=i}\\
&=2\pi i\times\left[\frac d{dz}\left\{(z-i)^2\,\frac{e^{i\omega z}}{(z-i)^2(z+i)^2}\right\}\right]_{z=i}\\
&=2\pi i\times\left[\frac{i\omega e^{i\omega z}(z+i)-2e^{i\omega z}}{(z+i)^3}\right]_{z=i}\\
&=2\pi i\times e^{-\omega}\left(\frac{i\omega(2i)-2}{(2i)^3}\right)\\
&=2\pi i\times e^{-\omega}\,\frac{\omega+1}{4i}=\frac\pi2(\omega+1)e^{-\omega}
\end{aligned}}
On the other hand
\EMdisp{\oint_C\frac{e^{i\omega z}}{(1+z^2)^2}\,dz\EMbr =\int_{-R}^{R}\frac{e^{i\omega x}}{(1+x^2)^2}\,dx+\int_0^\pi\frac{e^{i\omega Re^{it}}\,Rie^{it}}{(1+R^2e^{i2t})^2}\,dt}
and as \EMim{R\to\infty} the integral along the semicircle vanishes, while the imaginary part of the integral along the real axis (the odd function \EMim{\dfrac{\sin\omega x}{(1+x^2)^2}}) is zero:
\EMdisp{\oint_C\frac{e^{i\omega z}}{(1+z^2)^2}\,dz\ \xrightarrow{\ R\to\infty\ }\ 2\int_0^\infty\frac{\cos\omega x}{(1+x^2)^2}\,dx}
\EMdisp{\int_0^\infty\frac{\cos\omega x}{(1+x^2)^2}\,dx\EMbr =\frac\pi4(\omega+1)e^{-\omega}}

\end{EMexample}

\begin{EMunit}

\EMfigure{m14-semicircle}{The upper semicircular path; the double pole \EMim{z=i} lies inside the
path and \EMim{z=-i} outside it}

\end{EMunit}

\EMsection{Integrals of the Second Class}{Integrals of the Second Class}

\EMdisp{\int_0^{2\pi}f(a\cos t,\,b\sin t)\,dt}
To solve such integrals we put \EMim{z=e^{it}}. Then
\EMdisp{dz\EMbr =ie^{it}\,dt\;\EMbr \Longrightarrow\;dt\EMbr =\frac{dz}{ie^{it}}\EMbr =\frac{dz}{iz}}
\EMdisp{\sin t\EMbr =\frac{e^{it}-e^{-it}}{2i}\EMbr =\frac{z-z^{-1}}{2i},\qquad\EMbr  \cos t\EMbr =\frac{e^{it}+e^{-it}}2\EMbr =\frac{z+z^{-1}}2}

\begin{EMimportant}{}

\EMdisp{\int_0^{2\pi}f(a\cos t,\,b\sin t)\,dt\EMbr =\oint_Cf\!\left(a\,\frac{z+z^{-1}}{2},\;b\,\frac{z-z^{-1}}{2i}\right)\frac{dz}{iz}}

\end{EMimportant}

\begin{EMunit}

Here \EMim{C} is the unit circle, and the integral is easily solved with the method of residues.

\EMfigure{m14-unit-circle}{The unit circle \EMim{z=e^{it}}, \EMim{0\le t\le2\pi}}

\end{EMunit}

\begin{EMexample}{}

Evaluate the integral \EMim{\displaystyle\int_0^{2\pi}\frac1{2+\cos x}\,dx}.

\EMdisp{z\EMbr =e^{ix},\qquad\EMbr  dz\EMbr =ie^{ix}dx,\qquad\EMbr  dx\EMbr =\frac{dz}{iz},\qquad\EMbr  \cos x\EMbr =\frac{z+z^{-1}}2\EMbr =\frac{z^2+1}{2z}}
\EMdisp{2+\cos x\EMbr =2+\frac{z^2+1}{2z}\EMbr =\frac{z^2+4z+1}{2z}}
\EMdisp{\int_0^{2\pi}\frac1{2+\cos x}\,dx\EMbr =\oint_C\frac{2z}{z^2+4z+1}\,\frac{dz}{iz}\EMbr =\frac2i\oint_C\frac1{z^2+4z+1}\,dz\EMbr =\frac2i\times2\pi i\times\sum_k\operatorname{Res}\left\{\frac1{z^2+4z+1}\right\}\Bigg|_{z=z_k}}
\EMdisp{z^2+4z+1\EMbr =0\;\EMbr \Longrightarrow\;z_1\EMbr =-2+\sqrt3\EMbr \approx-0.27,\qquad\EMbr  z_2\EMbr =-2-\sqrt3\EMbr \approx-3.73}
Only \EMim{z_1} lies inside the unit circle.
\EMdisp{\int_0^{2\pi}\frac1{2+\cos x}\,dx\EMbr =4\pi\times\operatorname{Res}\left\{\frac1{z^2+4z+1}\right\}\Bigg|_{z=z_1}\EMbr =4\pi\left[\big(z+2-\sqrt3\big)\frac1{\big(z+2-\sqrt3\big)\big(z+2+\sqrt3\big)}\right]_{z_1=-2+\sqrt3}\EMbr =\frac{2\pi}{\sqrt3}}

\end{EMexample}

\begin{EMunit}

\EMfigure{m14-poles-cos}{The unit circle and the poles \EMim{z_1\approx-0.27} (inside) and
\EMim{z_2\approx-3.73} (outside)}

\end{EMunit}

\begin{EMexercise}{Applying the residue theorem to integrals}

Using the residue theorem, evaluate the following definite integrals:
\EMdisp{I_1\EMbr =\int_0^{2\pi}\frac{dt}{1-2p\cos t+p^2},\qquad\EMbr  I_2\EMbr =\int_{-\infty}^{\infty}\frac{\sin2x}{x^2+x+1}\,dx}
\EMdisp{I_3\EMbr =\int_{-\infty}^{\infty}\frac{x}{(x^2-2x+2)^2}\,dx,\qquad\EMbr  I_4\EMbr =\int_{-\infty}^{\infty}\frac{\cos^4x}{(x^2+1)(x^2+4)}\,dx}

\end{EMexercise}
